\section{Polyhedral tools and a feasible center}\label{sec:polyhedra}

The first step is to work with the exponentially many rank inequalities
without enumerating them.  We record the exact oracle algorithms we use,
then establish the geometric estimates needed to replace large coordinate
ranges by short ranges of part totals.  In particular, the distance estimate
below depends on the dimension, not on the numerical value of the rank.

Recall that for $x\in\mathbb R^E$ and $A\subseteq E$ we write
$x(A)=\sum_{j\in A}x_j$.  A set function $f:2^E\to\mathbb R$ is
\emph{normalized} if $f(\varnothing)=0$.  For any normalized submodular
function, including one that is not nondecreasing, put
\begin{equation}\label{eq:base-polytope}
 B(f)=\{x\in\mathbb R^E:x(E)=f(E),\quad x(A)\leq f(A)
                 \text{ for every }A\subseteq E\}.
\end{equation}
This convention lets us translate and reflect base polytopes.  The
singleton and complementary-singleton inequalities give
\[
 f(E)-f(E\setminus\{j\})\leq x_j\leq f(\{j\})
 \qquad (x\in B(f)).
\]
Thus $B(f)$ is bounded; for a nondecreasing $f$, its coordinates are
nonnegative.  Throughout, integrals in free coordinates use ordinary
Lebesgue measure with unit lattice mesh.  An integral in zero coordinates
means evaluation of its integrand.

\subsection{Exact oracle operations}

We use the following standard algorithmic results.  Their encoding
requirements are relevant because the capacity $R$ is given in binary.

\begin{proposition}[Oracle algorithms]\label{prop:oracle-tools}
The following operations have deterministic oracle and bit complexity
bounded by fixed polynomials in the indicated input sizes.
\begin{enumerate}
 \item An integer-valued submodular function on a finite ground set can
 be minimized, with a minimizing set returned, from an exact value oracle
 and a bound on the binary lengths of its values.  The size parameters are
 the ground-set cardinality and that value-length bound.
 \item Let $Q\subseteq\mathbb R^d$ be a bounded rational polytope, supplied
 by strong separation, a coordinate bound, and a bound on the binary
 encoding length of each individual defining inequality.  The separator
 returns inequalities of length polynomial in the query length and
 the supplied description bounds.  Exact rational
 linear optimization, including an emptiness decision and the production
 of an optimal vertex when $Q$ is nonempty, is polynomial in $d$, these
 bounds, and the objective length.  This assertion includes
 lower-dimensional polytopes.  All separation queries have polynomial
 encoding length.
\end{enumerate}
The first operation also applies to an integer submodular function with
a rational modular term added, and with prescribed elements required to
belong or not to belong to the minimizing set.
\end{proposition}

For the first assertion we use the submodular minimization theorem of
Iwata, Fleischer, and Fujishige \cite[Theorems 3.3 and 3.7]{IFF2001}.
For the second we use the strong separation--optimization equivalence
for well-described rational polyhedra
\cite[Definition 6.2.2, Theorem 6.4.9, and Remark 6.5.2]{GLS1988}.
The description size in that theorem bounds the encoding of one
inequality, and does not count the number of inequalities.  Its vertex
encoding bound also permits optimal-vertex selection by successive
rational face restrictions.
The LP bound counts calls to its separation oracle and additional bit
operations separately.  In our applications each such call is
implemented by submodular minimization and the explicit rational
arithmetic described next.

Here is how we supply these algorithms with the required oracles.
At a rational point $x$, separation from $B(r_i)$ minimizes
$r_i(A)-x(A)$ and checks the total equation.  Multiply this set function
by the product of the denominators of the coordinates of $x$.
The resulting integer submodular function has values of polynomial
binary length: the length of that product is at most the sum of the
denominator lengths.  Requiring inclusion of $F$ and exclusion of $G$,
where $F\cap G=\varnothing$, amounts to minimizing
$A\mapsto r_i(A\cup F)-x(A\cup F)$ on
$E\setminus(F\cup G)$.  Restrictions, contractions, and ranks on unions
of partition parts are handled in the same way.  All rational modular
translations used below have explicitly known denominators and
polynomial-length coefficients, so the same clearing of denominators
applies to them.  Additional coordinate bounds and part-total equations
are explicit linear constraints.

In particular, the original intersection
\begin{equation}\label{eq:original-polytope}
 P=B(r_1)\cap B(r_2)
\end{equation}
has a strong separation oracle and lies in $[0,R]^n$.  Its inequalities
have $0$--$1$ normals and right sides of $O(1+\log(R+1))$ bits.
Exact feasibility and optimization on $P$ therefore satisfy
Proposition~\ref{prop:oracle-tools} without any dependence polynomial
in the numerical value of $R$.

\subsection{Base-polytope calculus}

The following greedy identities, originating in the theory of
polymatroid polyhedra \cite{Edmonds1970}, identify the polytopes appearing
in translations, projections, and expansions later in the proof.

\begin{lemma}[Base-polytope identities]\label{lem:base-calculus}
Let $f,g:2^E\to\mathbb R$ be normalized submodular functions and let
$v\in\mathbb R^E$.
\begin{enumerate}
 \item If $w_{\pi(1)}\geq\cdots\geq w_{\pi(n)}$ and
 $A_k=\{\pi(1),\ldots,\pi(k)\}$, the vector
 \[
 b_{\pi(k)}=f(A_k)-f(A_{k-1})
 \]
 belongs to $B(f)$ and maximizes $w\cdot x$ there.
 \item With $(f+v)(A)=f(A)+v(A)$ and
 $f^{\mathrm{ref}}(A)=f(E\setminus A)-f(E)$, one has
 \[
 B(f+g)=B(f)+B(g),\qquad B(f+v)=B(f)+v,\qquad
 B(f^{\mathrm{ref}})=-B(f).
 \]
 \item For a partition $\mathcal D$ of $E$, define
 $\pi_{\mathcal D}(x)=(x(D):D\in\mathcal D)$ and
 $\bar f(\mathcal J)=f(\bigcup_{D\in\mathcal J}D)$.
 Then $\pi_{\mathcal D}(B(f))=B(\bar f)$.
 \item For distinct $a,b\in E$ and $t\geq0$, the segment
 $[-t,t](e_a-e_b)$ is $B(q_{ab})$, where
 \[
 q_{ab}(A)=t\,|\mathbf1_A(a)-\mathbf1_A(b)|.
 \]
 In particular, adding any collection of these segments to a base
 polytope again gives a base polytope, with integral rank if $f$ and
 the segment lengths are integral.
\end{enumerate}
\end{lemma}

\begin{proof}
When $E=\varnothing$, all assertions are immediate, so assume $n\geq1$.
Submodularity implies decreasing marginal increments: if $C\subseteq D$
and $j\notin D$, then
$f(C\cup\{j\})-f(C)\geq f(D\cup\{j\})-f(D)$.
Apply this inequality to the elements of any $A$ in their $\pi$-order.
The increments defining $b(A)$ are no greater than those along the
induced chain in $A$, and hence $b(A)\leq f(A)$.  Also $b(E)=f(E)$.
Summation by parts gives, for every $x\in B(f)$,
\[
 w\cdot x=w_{\pi(n)}f(E)
  +\sum_{k=1}^{n-1}(w_{\pi(k)}-w_{\pi(k+1)})x(A_k).
\]
All coefficients in the sum are nonnegative, and $b$ attains equality
in each of the resulting upper bounds.  This proves the first assertion
and, in particular, nonemptiness of $B(f)$.

The displayed support-function formula is additive in $f$.  Since
nonempty compact convex sets are determined by their support functions,
it proves the Minkowski-sum identity.  Translation follows directly from
the inequalities.  For reflection, the equation $x(E)=f(E)$ converts
$-x(A)\leq f(E\setminus A)-f(E)$ into the inequality on
$E\setminus A$ for $x$.  The reflected function is normalized and
submodular.

For the projection identity, containment in $B(\bar f)$ follows from
the inequalities on unions of parts.  Conversely, maximize a linear
functional on the projected coordinates.  Its lift to $\mathbb R^E$
is constant on each part.  In the greedy formula, place the parts in
decreasing order of these constants and place each part's elements
consecutively.  The increments telescope to the greedy formula for
$\bar f$.  The two projected sets thus have the same support function.

Finally, $q_{ab}$ is the cut function of a single edge and is
submodular, as is checked from its four possible membership patterns
on $\{a,b\}$.  Its base inequalities force every other coordinate to
be zero, force $x_a+x_b=0$, and give $|x_a|\leq t$.  This proves the
last assertion.
\end{proof}

The next construction is the expansion of an integral polymatroid to
its natural matroid \cite{Helgason1974}.  We use it only to prove
inequalities for weights.  Its ground set may have size comparable with
$R$, and we never construct that ground set in the algorithm.

\begin{lemma}[The expansion matroid]\label{lem:expansion-matroid}
Let $r:2^E\to\mathbb Z_{\geq0}$ be normalized, nondecreasing, and
submodular.  Replace each $j\in E$ by any finite number of labelled
copies.  Declare a set of copies independent when its occupancy vector
$z$ satisfies $z(A)\leq r(A)$ for every $A\subseteq E$.
This defines a matroid.  A set of copies with occupancy vector $v$ has
matroid rank
\begin{equation}\label{eq:expansion-rank}
 m(v)=\min_{A\subseteq E}\{r(A)+v(E\setminus A)\}.
\end{equation}
\end{lemma}

\begin{proof}
Hereditary closure is immediate.  Let $I,J$ be independent copy sets
with occupancy vectors $z,z'$ and $|I|<|J|$.  If no element of $J\setminus I$
can be added to $I$, then, for every coordinate $j$ with $z'_j>z_j$,
some set $A_j$ containing $j$ is tight for $z$.  Indeed, failure of the
increment $z+e_j$ yields such a set because both its slack and its rank
are integral.  Tight sets are closed under union and intersection:
submodularity and modularity give
\[
 r(A)+r(B)=z(A)+z(B)=z(A\cup B)+z(A\cap B)
 \leq r(A\cup B)+r(A\cap B)\leq r(A)+r(B).
\]
Their union $A$ is therefore tight.  On $E\setminus A$ we have
$z'_j\leq z_j$, whereas $z'(A)\leq r(A)=z(A)$.  This contradicts
$|J|>|I|$ and proves augmentation.

Every independent subprofile $z\leq v$ has
$z(E)\leq r(A)+v(E\setminus A)$ for every $A$.
Choose a maximal independent subprofile of $v$.  Each coordinate where
$z_j<v_j$ is blocked by a tight set containing $j$.  The union $A$ of
these sets is tight, and $z=v$ outside $A$.  Thus
$z(E)=r(A)+v(E\setminus A)$, proving equality in
\eqref{eq:expansion-rank}.  If no coordinate is blocked, take
$A=\varnothing$.
\end{proof}

\subsection{Integrality and quantitative stability}

We require two kinds of systems: the intersection of two base
polytopes, and a single base polytope with prescribed totals on the
parts of a partition.  These systems share a useful structure at every
feasible point.  The structure persists after coordinate bounds are
added, including bounds defining a cube.  We give the uncrossing proof
of the classical integrality assertion \cite{Edmonds1970} together with
the distance and slope bounds that we need.

\begin{lemma}[Polyhedral controls]\label{lem:polyhedral-controls}
Let $Q\subseteq\mathbb R^E$ be nonempty and be defined either by
$B(f_1)\cap B(f_2)$ or by $B(f)$ together with equations
$x(D)=a_D$ for the parts $D$ of a partition of $E$.
The functions are normalized and submodular.  In either case arbitrary
coordinate lower and upper bounds may also be imposed.
\begin{enumerate}
 \item If the ranks and all other right sides are integral, every
 vertex of $Q$ is integral.
 \item Suppose $x\in\mathbb R^E$ satisfies every defining inequality
 with additive violation at most $b\geq0$, and every defining equality
 with absolute error at most $b$.  Then
 \begin{equation}\label{eq:polyhedral-distance}
  \operatorname{dist}_{\infty}(x,Q)\leq10(n+1)^2b.
 \end{equation}
 \item Fix $a\in Q$ and an affine function $F$ whose linear part is
 integral.  The function
 \[
  \phi(l)=\max\{F(y):y\in Q,\ \|y-a\|_\infty\leq l\},
  \qquad l\geq0,
 \]
 is continuous, nondecreasing, concave, and piecewise affine.  Every
 positive slope of an affine piece is an integer and hence at least
 one.  Eventually $\phi(l)=\max_Q F$.
 Consequently, if $F\geq0$ on $Q$ and $\phi(L)<L$ for some $L>0$,
 then $\phi(L)=\max_QF$.
\end{enumerate}
\end{lemma}

\begin{proof}
We first establish the structure of the tight constraint normals.
For a fixed rank system at $y\in B(f)$, the tight sets are closed
under union and intersection by the calculation in the preceding
proof.  Choose a maximal strictly increasing chain of tight sets from
$\varnothing$ to $E$.  No tight set can cut a difference between
successive members of this chain: otherwise intersecting that set with
the upper member and adjoining the lower member inserts another tight
set.  Every tight-set indicator is therefore a sum of these difference
indicators.  In particular, the chain indicators span all tight rank
normals.

A second version preserves nonnegative coefficients.  Any nonnegative
combination of tight-set indicators can be expressed as a nonnegative
combination of indicators from a chain.  To see this, keep the resulting
vector and the total coefficient fixed, allowing a coefficient on the
empty set, and maximize the sum of the coefficients times $|A|^2$.
The feasible coefficient set is compact.  Two incomparable sets with
positive coefficients can be replaced, with the smaller coefficient,
by their union and intersection.  The vector and total coefficient
are unchanged, while the maximized quantity increases by
$2|A\setminus B|\,|B\setminus A|$ times that coefficient.
Thus no such pair remains.  Signed multiples of the total equation
can be kept separately.

The matrix formed by the indicators of two chains is totally
unimodular.  More generally this holds for two laminar families,
and therefore for a chain and a partition.  We apply Ghouila-Houri's
row-signing criterion
\cite[Chapter II, Section 3, fundamental theorem]{GH1964}.
For any selection of rows from the first
family, alternate signs along nested selected sets, starting with
$+1$ on the maximal sets.  In any column their sum is either $0$ or
$1$.  Start with $-1$ in the second family, obtaining $0$ or $-1$.
Their combined sum is in $\{-1,0,1\}$.  Coordinate rows may be added:
choose their signs separately in each column to retain this bound.
Negating rows or adding signed copies also preserves total
unimodularity.  Thus the same conclusion holds with either orientation
of equality normals and with coordinate bounds.

At a vertex, the tight normals span $\mathbb R^E$.  Replace the tight
rank rows by their chain spans, retain the partition and coordinate
rows, and choose $n$ independent rows.  Their square matrix is
unimodular.  With integral right sides its unique solution is integral,
proving the first assertion.

For the distance estimate, let $y\in Q$ minimize $\|x-y\|_\infty$
and write $t=\|x-y\|_\infty$.  If $t=0$ there is nothing to prove.
The optimality condition for this convex minimization supplies a
normal $g$ to $Q$ at $y$ satisfying
\[
 \|g\|_1=1,\qquad g\cdot(x-y)=t.
\]
Express $g$ as a nonnegative combination of the active inequality
normals and the two orientations of active equality normals.
Uncross the rank contributions separately as above.  The remaining
normals belong to a totally unimodular matrix.  A conic representation
can be chosen with linearly independent normals, say
$g=\sum_{k=1}^m\mu_k a_k$ with $m\leq n$ and $\mu_k\geq0$.
Choose $m$ columns giving a nonsingular square minor of the row
matrix.  That minor has determinant $\pm1$ and every cofactor is
$0$ or $\pm1$.  Solving for the coefficients gives
$\mu_k\leq\|g\|_1=1$.  Each active oriented row satisfies
$a_k\cdot(x-y)\leq b$, so
\[
 t=\sum_{k=1}^m\mu_k a_k\cdot(x-y)\leq nb.
\]
This stronger estimate implies \eqref{eq:polyhedral-distance}.

For the last assertion, the cube constraint is
$y_j\leq a_j+l$ and $-y_j\leq-a_j+l$.
Every vertex of the constrained system has a basis selected by the
same chain-span argument.  Its basis matrix is unimodular and its
right side has the form $b_0+l b_1$ with $b_1$ integral.  Thus that
vertex depends affinely on $l$, with integral derivative wherever
the basis is feasible.  There are only finitely many possible bases.
Optimizing $F$ consequently gives a piecewise affine function whose
slopes are integers.  Inclusion of cubes gives monotonicity; convex
combinations of feasible points for radii $l_1,l_2$ give concavity.
Continuity on $(0,\infty)$ follows from finite concavity, and at zero
from $|F(y)-F(a)|\leq\|\nabla F\|_1l$.  Boundedness of $Q$ gives
eventual equality with its unrestricted maximum.

If $\phi(L)<\max_QF$, concavity and monotonicity exclude a
zero-slope interval before $L$.  Every preceding affine piece then
has slope at least one, so
$\phi(L)\geq\phi(0)+L\geq L$ when $F\geq0$ on $Q$.
This proves the stated consequence.
\end{proof}

\subsection{Constructing a feasible center}

The distance estimate controls the effect of changing part totals.
We also need one feasible point at which a small slack certifies that
the same slack is small throughout the feasible polytope.  The following
construction works for any bounded rational polytope supplied by the
oracle interface, including the lower-dimensional slices used later.
It is the affine-simplex version of the determinant-replacement
construction for barycentric spanners
\cite[Section 2.3 and Proposition 2.4]{AK2004}; we include the argument
to obtain the precise slack and encoding bounds.

\begin{lemma}[A center for affine slacks]\label{lem:center}
Let $Q\subseteq\mathbb R^n$ be a nonempty bounded rational polytope
satisfying the hypotheses of Proposition~\ref{prop:oracle-tools}.
One can find a rational $c\in Q$, with polynomial encoding length,
such that every affine $F$ nonnegative on $Q$ satisfies
\begin{equation}\label{eq:center-slack}
 \max_{x\in Q}F(x)\leq4(n+1)F(c).
\end{equation}
The construction has polynomial oracle and bit complexity.
It also returns an affinely spanning vertex simplex of $Q$ such that
every point of $Q$ has all barycentric coordinates of absolute value
at most two with respect to that simplex.
\end{lemma}

\begin{proof}
Start from a vertex $v_0$ obtained by linear optimization.
Given vertices whose affine span has direction space $L$, optimize
both signs of a rational basis of $L^\perp$.  Either every optimum
equals its value at $v_0$, proving that $Q$ lies in the current span,
or an optimal vertex extends the span.  After at most $n$ extensions
this gives $d+1$ affinely independent vertices, where $d=\dim Q$.
For $d=0$, take $c=v_0$ and stop.

Choose a rational coordinate projection that is one-to-one on the
affine hull.  For the current simplex $v_0,\ldots,v_d$, let
$\lambda_0,\ldots,\lambda_d$ be its affine barycentric-coordinate
functions on that hull.  Optimize both signs of each $\lambda_j$
over $Q$.  If some value has absolute value greater than two, replace
$v_j$ by an attaining vertex.  The absolute determinant of the
simplex in the fixed coordinate projection is multiplied by
$|\lambda_j|>2$.

All vertices of a well-described rational polytope have polynomial
encoding length \cite[Lemma 6.2.4]{GLS1988}.  Consequently every
nonzero simplex determinant encountered lies between
$2^{-B}$ and $2^B$ for a polynomially bounded $B$.  There are at most
$2B+1$ replacements.  Determinants, barycentric functions, objectives,
and returned vertices throughout the procedure have polynomial
encoding length, so the entire construction has polynomial cost.

At termination, $|\lambda_j(x)|\leq2$ for every $x\in Q$ and every
$j$.  Set $c=(d+1)^{-1}\sum_{j=0}^d v_j$.  If $F\geq0$ on $Q$,
then $F(v_j)\geq0$, and for every $x\in Q$,
\[
 F(x)=\sum_{j=0}^d\lambda_j(x)F(v_j)
 \leq2\sum_{j=0}^dF(v_j)=2(d+1)F(c).
\]
This implies \eqref{eq:center-slack}.
\end{proof}

We can now make the preliminary branches of the counting algorithm
exact.  If $R=0$, return $1$; the only feasible vector is zero,
including when $E=\varnothing$.  Otherwise use
Proposition~\ref{prop:oracle-tools} to test $P$ in
\eqref{eq:original-polytope}.  If it is empty, return $0$.
If it is nonempty, Lemma~\ref{lem:polyhedral-controls} gives an
integral vertex, so $Z\geq1$.  Obtain a rational center $c\in P$
from Lemma~\ref{lem:center}.  Henceforth this nonempty case is
understood.  In particular the algorithm will return zero on every
execution for every instance with $Z=0$.
